%% Reciprocal proofs and exact source sketch boundaries for Lawvere1970.

\begin{annexitem}{annex:p10.lawvere-sheaves.internal}{thm:p10-lawvere-sheaves-internal-adjoints}{The original internal-completeness announcement}
\begin{proof}
Theorem~\ref{thm:p10-lawvere-sheaves-dependent-product} constructs
\(\Pi_f\) from power-object functional graphs, gives its counit
and both inverse transpositions, and proves the adjunction for all
families and all generalized pullback bases.  It also constructs
the complete original package in each slice.  Theorem
\ref{thm:p10-lawvere-sheaves-predicate-adjoints} then constructs
internal Heyting operations, both predicate adjoints and stable
images from those adjunctions, without assumed exactness or
Boolean negation.  Theorem
\ref{thm:p10-lawvere-sheaves-effective-quotients} proves regularity
and effectiveness of internal equivalence relations from the
constructed images.

Theorem~\ref{thm:p10-lawvere-sheaves-internal-presheaf-topos}
constructs finite limits and finite colimits of actions, compatible
exponentials with both inverse transpositions, and the internal
sieve classifier with unique characteristic pullbacks.  Thus it
proves that \(C^{\mathrm{op}}\mathcal E\) is a topos for
every internal category \(C\).  Theorem
\ref{thm:p10-lawvere-sheaves-internal-kan} constructs the right
extension by compatible dependent assignments and the left
extension by an explicit finite parallel-pair coequalizer, and
proves both adjunctions by natural inverse correspondences.
These results give every structural and adjoint assertion of the
body theorem at the original internal scope.

All prerequisites in these deductions are constructed from the
original finite-limit, finite-colimit, exponential and classifier
package; the dependency order starts with functional graphs,
not with the desired adjoint or exactness conclusions.  The
one-object internal group case is an action category, and a
discrete internal category on \(X\) gives \(\mathcal E/X\),
so the construction genuinely includes both source regimes.
Its products and compatibility conditions are internal families,
not arbitrary externally indexed limits.  An externally countable
topos is therefore not excluded by this form of completeness.
\end{proof}

These are editorial full deductions.  Lawvere's p.~330 supplies the
announcement and dependency program, not the detailed constructions
above; their presence here does not retrospectively promote the
article's historical proof standing.
\end{annexitem}

\begin{annexitem}{annex:p10.lawvere-sheaves.dependent-product}{thm:p10-lawvere-sheaves-dependent-product}{Functional relations, evaluation, and all pullback bases}
\begin{proof}
\paragraph{The constructions before evaluation.}
The power-object correspondence follows by composing the given
exponential and classifier bijections.  It is natural because
transposition is natural and classification commutes with pullback.
The conjunction map classifies the pullback intersection of the
two subobjects \(\mathsf{true}\times T\) and
\(T\times\mathsf{true}\) of \(T\times T\).  For any
\(u,v:Z\to T\), its composite with \((u,v)\) therefore
classifies the intersection of their truth pullbacks.  In particular
\(u\land v=u\) means inclusion of those subobjects.  This
uses neither implication nor a universal quantifier.

The two graph conditions define actual subobjects by transposing
the corresponding maps into \(T^{X\times E}\) and
\(T^{X\times E\times E}\), taking their equalizers, and
intersecting the equalizer subobjects of \(D_0\).
Consequently they hold after every map of parameters into \(D_1\),
not just at global points.  The restricted relation \(R_1\)
is the truth pullback of its membership predicate.

To prove that \(d:R_1\to D_1\times X\) is monic, take
\(s,t:W\to R_1\) with \(ds=dt\).  Their \(D_1\) and
\(X\) coordinates agree; call their \(E\) coordinates
\(e,e':W\to E\).  Both membership predicates are true on
\(W\).  The second graph equation therefore says
\(\delta_E(e,e')=\mathsf{true}\,!_W\).  Classification of
the diagonal gives \(e=e'\).  All their product coordinates
agree, and \(R_1\) is a subobject of that product, so \(s=t\).
This proves monicity for arbitrary \(W\).  The first equation
likewise gives \(p(e)=x\) and \(f(x)=y\) on the relation,
so \(d\) factors through \(D_1\times_Y X\).

Now pull this mono back to \(D\).  By the equalizer defining
\(D\), its classifier is exactly \(\gamma|_{D\times X}\).
The latter also classifies \(D\times_Y X\hookrightarrow
D\times X\).  Two monos with this same classifier are
canonically isomorphic over their codomain, by the pullback
universal property.  We thus obtain
\[
                       R_D\cong D\times_Y X.
\]
Projecting the relation to its \(E\) coordinate gives the
displayed counit \(\varepsilon_p\).  The first graph equation
gives \(p\varepsilon_p=\operatorname{pr}_X\).
No epi-mono image factorization, Boolean complement, or inference
from an internally inhabited relation to a chosen section appears
in this step: totality is equality of the actual monomorphisms.

\paragraph{The inverse transpositions.}
Let \(a:Z\to Y\), put \(G_a=Z\times_Y X\), and take
\(h:G_a\to E\) with \(ph=\operatorname{pr}_X\).
Its graph is the mono
\[
 G_a\longrightarrow Z\times X\times E,
       \qquad (z,x)\longmapsto(z,x,h(z,x)).
\]
It is monic because its first two coordinates are the mono
\(G_a\hookrightarrow Z\times X\).
Power-object classification transposes its predicate to a unique
\(r_h:Z\to T^{X\times E}\).  The pair \((a,r_h)\)
gives a map into \(D_0\).  The graph lies over \(p\) and
\(f\), and is single-valued; these facts make it satisfy both
defining equalities of \(D_1\), hence factor through \(D_1\).
For clarity, an inclusion of relation subobjects means precisely
an equation of their conjunction classifiers, and equality of
two output coordinates means the truth pullback of the diagonal;
these are exactly the two equalizer equations, not an assumed
internal proof calculus.

The pullback of \(R_1\) along this parameter map is the graph
just classified.  Its projection to \(Z\times X\) is
isomorphic to \(G_a\hookrightarrow Z\times X\).
Classifier pullback naturality makes the two transposed domain
maps equal after \(Z\to D_1\).  The parameter map therefore
factors uniquely through \(D\), giving \(\Lambda_f h\).

Conversely a map \(k:Z\to D\) over \(Y\) pulls the counit
back to a map
\[
 h_k:G_a\longrightarrow E.
\]
The relation classified by its \(T^{X\times E}\) coordinate
is the pullback of \(R_D\).  The isomorphism
\(R_D\cong D\times_Y X\) identifies it with the graph of
\(h_k\), including its output coordinate.  Uniqueness of
classification recovers that coordinate of \(k\); its \(Y\)
coordinate is already \(a\).  Since \(D\hookrightarrow D_0\)
is monic, \(\Lambda_fh_k=k\).  Starting instead with \(h\),
the graph identification immediately gives
\(h_{\Lambda_fh}=h\).  These prove both inverse identities
and uniqueness of transposition.

For \(z:Z'\to Z\) over \(Y\), graph pullback gives
\[
 \Lambda_f\bigl(h\circ(z\times_Y1_X)\bigr)
                                  =(\Lambda_fh)z.
\]
For \(v:E\to E'\) over \(X\), define
\(\Pi_fv\) by transposing \(v\varepsilon_p\).
The two inverse identities prove
\(\Pi_f1=1\) and \(\Pi_f(v'v)=(\Pi_fv')(\Pi_fv)\).
They also prove
\[
              \Lambda_f(vh)=(\Pi_fv)(\Lambda_fh).
\]
Thus the constructed assignment is a functor and the bijection
is natural in both variables.  The unit is the transpose of the
identity of \(f^*Z\); the two proved inverse identities are
the unit/counit triangle identities.  This establishes the
adjunction, rather than merely a family of candidate objects.

\paragraph{Arbitrary generalized pullback bases.}
In the stated pullback square, fix any \(z:Z\to Y'\).
The pullback universal property gives a canonical isomorphism
\[
                Z\times_{Y'}X'\cong Z\times_Y X,
\]
where \(Z\to Y\) is \(bz\).  A map from the left-hand
object to \(a^*E=X'\times_XE\) over \(X'\) is exactly
a map from the right-hand object to \(E\) over \(X\).
The already proved transposition bijections give, naturally in
every such \(Z\),
\[
\begin{aligned}
 \mathcal E/Y'(Z,b^*\Pi_fp)
 &\cong\mathcal E/Y(Z,\Pi_fp)\\
 &\cong\mathcal E/X(Z\times_YX,E)\\
 &\cong\mathcal E/X'(Z\times_{Y'}X',a^*E)\\
 &\cong\mathcal E/Y'(Z,\Pi_{f'}a^*p).
\end{aligned}
\]
In the first line the middle hom collection means maps over the
fixed map \(bz:Z\to Y\), as prescribed by the pullback.
The representers are canonically isomorphic.  More explicitly,
pull back \(\varepsilon_p\) along \(b\), use the displayed
domain isomorphism, and transpose along \(f'\); this constructs
the usual comparison map.  The above bijections show that it is
an isomorphism, with an inverse corresponding to the reverse
domain isomorphism.  They include all generalized \(Z\), even
initial or noninhabited ones, not only global points.

For an identity square the comparison has the original counit,
so uniqueness of transposition makes it the identity.  For two
pasted pullback squares, both comparisons have the same twice
pulled-back counit under the associativity isomorphism of
pullbacks.  The same uniqueness equates the two composites.
For composable \(f:X\to Y\), \(g:Y\to V\), successive
transposition gives
\[
 \mathcal E/V(Z,\Pi_g\Pi_fE)
 \cong\mathcal E/Y(g^*Z,\Pi_fE)
 \cong\mathcal E/X(f^*g^*Z,E).
\]
Since \(f^*g^*Z\cong(gf)^*Z\), this proves
\(\Pi_{gf}\cong\Pi_g\Pi_f\); the identity case proves
\(\Pi_{1_X}\cong1\).

\paragraph{The full slice package.}
Finite limits in \(\mathcal E/X\) are constructed from the
terminal object \(1_X:X\to X\), products as fiber products
over \(X\), and base equalizers with their inherited map to
\(X\).  These constructions generate all finite limits;
the empty limit is \(1_X\), not the base terminal object
with an unspecified map to \(X\).  Finite colimits are the base colimits with the
structure map induced by the compatible maps to \(X\):
the base universal cocone gives exactly the slice universal
cocone.  This includes \(0\to X\) as initial object.

The classifier in the slice is \(X\times T\to X\),
with truth \((1_X,\mathsf{true}\,!_X):X\to X\times T\).
A mono \(A\hookrightarrow B\) over \(X\) is also a
base mono: for base maps into \(A\) with equal composites,
their induced maps to \(X\) coincide, so slice cancellation
applies.  The reverse implication is immediate.  Its classifier
in the slice is the pair of \(B\to X\) with the unique
base classifier \(B\to T\).  Base pullback classification
gives the slice pullback and uniqueness.

For \(p:E\to X\), \(q:F\to X\), put
\(H=\Pi_p(p^*q)\).  For any \(Z\to X\), the adjunction
just proved identifies maps \(Z\to H\) over \(X\) with
maps \(Z\times_XE\to E\times_XF\) over \(E\).
The latter are precisely maps \(Z\times_XE\to F\) over
\(X\).  This proves the exponential universal property,
including its evaluation and inverse transposition identities.
All four original structural assumptions hold in the slice.
The construction never imposed nonempty fibers, choice, classical
logic, or an external infinite limit.
\end{proof}
\end{annexitem}

\begin{annexitem}{annex:p10.lawvere-sheaves.predicate-adjoints}{thm:p10-lawvere-sheaves-predicate-adjoints}{Non-Boolean quantification and the actual image construction}
\begin{proof}
\paragraph{Universal quantification first.}
Applying \(\Pi_f\) to a mono into the terminal family of
\(\mathcal E/X\) gives a mono into the terminal family
of \(\mathcal E/Y\), since a right adjoint preserves
terminal objects and monomorphisms.  Define \(\forall_fA\)
to be this subobject.  For \(B\hookrightarrow Y\), the
slice adjunction gives
\[
        B\leq\forall_fA\quad\Longleftrightarrow\quad
                                           f^*B\leq A.
\]
The arbitrary-pullback comparison for \(\Pi\) gives the
same comparison for these monos.  It identifies subobjects,
so it is equality of classified subobjects, proving
\(b^*\forall_fA=\forall_{f'}a^*A\).
No formula \(\neg\exists\neg\) was used.

\paragraph{An image constructed without assumed exactness.}
Let \(v:E\to X\) be arbitrary.  Let \(K\hookrightarrow T^X\)
be the equalizer of precomposition \(T^v:T^X\to T^E\)
and the constant top predicate \(\mathsf{true}_E:1\to T^E\).
Thus its predicates are exactly those containing \(v\).
Restrict universal membership to \(K\times X\) and,
after interchanging factors, call it
\(R\hookrightarrow X\times K\).  For
\(\pi:X\times K\to X\), put
\[
                         I=\forall_\pi R\hookrightarrow X.
\]
This is a defined object, using the already constructed \(\Pi\).
The equality defining \(K\) says that pulling \(R\) back
along \(v\times1_K:E\times K\to X\times K\) gives
the whole object \(E\times K\).  Universal base change
therefore gives \(v^*I=\forall_{E\times K\to E}\mathsf{top}
=\mathsf{top}\).  Classification factors \(v\) uniquely
through \(i:I\hookrightarrow X\); write \(v=iq\).

If \(B\hookrightarrow X\) also contains \(v\), its
classifier transposes to a point \(b:1\to K\), by the
defining equalizer.  The universal-quantifier counit says
\(\pi^*I\leq R\).  Pulling this inclusion back along
\(1_X\times b:X\to X\times K\) gives \(I\leq B\).
Hence \(I\) is the least subobject through which \(v\)
factors, not a name for an unconstructed image.

To prove \(q:E\to I\) epic, take maps \(s,t:I\to W\)
with \(sq=tq\), and let \(L\hookrightarrow I\) be their
equalizer.  The map \(v\) factors through
\(L\hookrightarrow I\hookrightarrow X\).  The just-proved
minimality implies \(I\leq L\) as subobjects of \(X\).
Thus \(L=I\) and \(s=t\).  This is cancellation for every
\(W\), so \(q\) is epic.

Any epi--mono factorization \(v=m e\) has that same minimal
subobject.  Indeed, if a subobject \(B\) contains \(v\),
its classifier satisfies \(\chi_Bme=\mathsf{true}\,!e\).
Cancel the epi \(e\) to obtain
\(\chi_Bm=\mathsf{true}\,!\), so \(m\) factors through
\(B\).  This proves uniqueness of the image subobject.
It also shows that a mono which is epic is an isomorphism:
its classifier is top by epi cancellation, and its classifier
pullback is therefore the identity subobject.

\paragraph{Stability does not assume regularity.}
A base epi over \(B\) is epic in \(\mathcal E/B\), by
the same cancellation.  Conversely an epi in that slice is
a base epi: two ordinary maps to \(W\) can be paired with
their common structure map to give maps into \(B\times W\)
in the slice, where cancellation applies.
Pullback \(a^*:\mathcal E/B\to\mathcal E/A\) has the
right adjoint just constructed.  A left adjoint preserves
epimorphisms: if two maps agree after its image of an epi,
their adjoint transposes agree after that epi, hence agree;
inverse transposition cancels the image.  Consequently every
pullback of a base epi is a base epi.  Pulling an image
factorization back gives an epi followed by a mono.  Its
uniqueness just proved makes it the pulled-back map's image.
This proves stable images from the original package without
assuming the desired exactness theorem as a premise.

\paragraph{Existential quantification and the pullback identities.}
Define \(\exists_fA=\operatorname{Im}(A\to X\to Y)\).
Its minimality proves
\(\exists_fA\leq B\) iff \(A\to Y\) factors through
\(B\), equivalently \(A\leq f^*B\).  Stability of images
in the displayed pullback square proves the existential
base-change equality.  Pulling the image of \(A\to Y\)
back to \(B\hookrightarrow Y\) gives exactly the image
of \(A\times_YB\to B\).  Regarded as subobjects of
\(Y\), this states
\(\exists_f(A\cap f^*B)=\exists_fA\cap B\).
No merely commuting square was substituted for a pullback.

\paragraph{Heyting operations.}
The initial object is strict: \((-)\times Z\), being a left
adjoint, preserves \(0\).  A map \(Z\to0\) gives a
section of \(Z\times0\to Z\), so \(Z\) is a retract
of an initial object and is initial.  It follows that
\(0\to X\) is monic, giving bottom; \(1_X\) gives top.
Intersections give meets.  The image of
\(A+B\to X\) is the least subobject containing both
\(A\) and \(B\), by the coproduct universal property and
image minimality, so it is their join.

Let \(L\hookrightarrow T\times T\) be the equalizer
of \(\land\) and \(\operatorname{pr}_1\), and let
\(\Rightarrow:T\times T\to T\) classify it.  For
\(A,B\hookrightarrow X\), classify their implication
by \(\Rightarrow(\chi_A,\chi_B)\).  A mono
\(D\hookrightarrow X\) lies in this truth pullback iff
the restricted classifiers satisfy
\[
                (\chi_A\land\chi_B)|_D=\chi_A|_D,
\]
iff \(D\cap A\leq B\).  This proves the Heyting
residuation with no complement reasoning.  In particular
intersection with \(A\), having this right adjoint,
preserves the binary joins and bottom just constructed:
an upper-bound comparison through residuation proves
both inequalities of the distributive law.  The subobject
lattices are Heyting algebras.  Pulling back the displayed
classifier formulas preserves top, bottom, intersections,
and implication; stable image plus preservation of finite
coproducts by pullback (its adjoint pair was constructed)
preserves joins as well.

Construct the corresponding maps on \(T\) using universal
predicates on \(T\times T\) and the just-defined subobject
operations.  Their algebra equations hold after every
generalized predicate, and in particular after the universal
ones on finite products of \(T\), so they are actual
morphism equations.  The residuation relation is the same
equalizer relation \(L\); hence this is an internal
Heyting algebra, not just a collection of global Boolean
algebras.

\paragraph{Internal rather than merely external adjoints.}
Let \(U_X\hookrightarrow T^X\times X\) be universal
membership and let \(s=1_{T^X}\times f\).  The mono
\(\Pi_sU_X\hookrightarrow T^X\times Y\) has a
classifier whose transpose is a map
\(\forall_f:T^X\to T^Y\).  The image of
\(U_X\to T^X\times Y\) similarly transposes to
\(\exists_f:T^X\to T^Y\).
For every \(r:Z\to T^X\), universal membership pulls
back to its classified relation \(R\subseteq Z\times X\).
The two actual pullback-stability results therefore identify
\(\forall_fr\) and \(\exists_fr\) with the operations
over \(1_Z\times f\) on \(R\).  Their order adjunctions
hold in every such parameter context by the already proved
subobject adjunctions.  This proves the internal order
adjunction identities, including generalized predicates and
all further parameter pullbacks.
Right adjoints preserve top, and images of \(0\) remain
\(0\); no added inhabited-sort assumption is used.
\end{proof}
\end{annexitem}

\begin{annexitem}{annex:p10.lawvere-sheaves.effective-quotients}{thm:p10-lawvere-sheaves-effective-quotients}{Descent along epis and quotients by equality of rows}
\begin{proof}
\paragraph{Every epi coequalizes its kernel pair universally.}
Let \(q:X\twoheadrightarrow Y\) and suppose
\(h:X\to Z\) has equal composites with the two projections
of \(X\times_YX\).  Form the already constructed image
of \((q,h):X\to Y\times Z\):
\[
           X\xrightarrow{e}\Gamma\hookrightarrow Y\times Z,
\]
where \(e\) is epic.  Let \(d:\Gamma\to Y\) and
\(z:\Gamma\to Z\) be the coordinate maps.  Since
\(q=de\) is epic, \(d\) is epic.

The map
\[
             X\times_YX\longrightarrow\Gamma\times_Y\Gamma
\]
induced by \(e\) in both coordinates is epic: it factors
as two pullbacks of \(e\), whose epimorphy was established
in the stable-image proof.  The two \(Z\) coordinates on
\(\Gamma\times_Y\Gamma\) agree after this epi by the
hypothesis on \(h\), and hence agree.  Their \(Y\)
coordinates agree by the fiber product.  Monicity of
\(\Gamma\hookrightarrow Y\times Z\) therefore equates
the two projections of \(\Gamma\times_Y\Gamma\).
For any two maps into \(\Gamma\) equalized by \(d\),
the fiber-product universal property now equates them;
thus \(d\) is monic.  The proved mono-epi criterion makes
\(d\) an isomorphism.  Put \(\bar h=z d^{-1}:Y\to Z\).
Then \(\bar hq=h\), and this factorization is unique
because \(q\) is epic.  This is exactly the kernel-pair
coequalizer universal property.  Together with the
previously proved stable image factorizations it establishes
regularity; it was not a premise of their construction.

\paragraph{An arbitrary internal equivalence relation.}
Let \(\rho:X\times X\to T\) classify \(R\), and let
\(r:X\to T^X\) be its transpose.  If \((x,y)\) belongs
to \(R\), symmetry and transitivity give
\[
                         \rho(x,z)=\rho(y,z)
\]
for every generalized \(z\).  More explicitly, over the
parameter object \(R\), the two membership subobjects of
\(R\times X\) have inclusions in both directions given
by composing the equivalence relation with the fixed pair
and its symmetric pair.  Their classifiers are equal;
exponential transposition gives equality of the two row
maps on \(R\).  Hence \(R\) lies in the kernel pair
of \(r\).

Conversely, on that kernel pair the row maps are equal.
Evaluate them at the first coordinate \(x\).  Reflexivity
gives \(\rho(x,x)=\mathsf{true}\), so row equality
gives \(\rho(y,x)=\mathsf{true}\); symmetry puts
\((x,y)\) in \(R\).  The classifier pullback property
therefore factors the kernel-pair mono through \(R\).
The two inclusions prove equality of subobjects, without
Boolean complement or choice of an element of each class.

Take the image factorization \(r=mq\), with \(m\)
monic and \(q:X\twoheadrightarrow Q\).  Since \(m\)
is monic, the kernel pairs of \(r\) and \(q\) coincide,
so that of \(q\) is \(R\).  The first paragraph makes
\(q\) the coequalizer of \(R\); every map equating
\(R\) factors through it uniquely.  This constructs an
effective quotient of every internal equivalence relation.
It also identifies it with any other coequalizer of that
relation by the two universal properties.  The construction
and its proof use the already supplied finite limits,
power objects, dependent products, and stable images.
No exactness theorem or externally accumulated quotient
collection is assumed.
\end{proof}
\end{annexitem}

\begin{annexitem}{annex:p10.lawvere-sheaves.internal-presheaf-topos}{thm:p10-lawvere-sheaves-internal-presheaf-topos}{Actions, compatible exponentials, and internal sieves}
\begin{proof}
\paragraph{The free and cofree family actions.}
The displayed underlying families exist by finite pullbacks,
composition \(\Sigma_s\), and the constructed \(\Pi_t\).
An element of \(L_CB\) is \((a:d\to c,b\in B_c)\),
and restriction along \(u:d'\to d\) is \((a u,b)\).
The category identities and associativity prove the presheaf laws.
A family map \(k:B\to U_CE\) gives an action map
\[
              L_CB\to E,\qquad(a,b)\mapsto k(b)\cdot a.
\]
Conversely an action map gives \(k(b)\) by evaluation at
\((i(c),b)\).  An action map must have the displayed values,
since \((a,b)=(i(c),b)\cdot a\).  These two assignments
are inverse and natural, proving \(L_C\dashv U_C\).
In map notation the displayed formula is a composite of the
pullback of \(k\) with the action, so no global-point
assumption is involved.

For \(R_CB\), an assignment at \(c\) has \(b_a\in B_d\)
for each \(a:d\to c\).  The counit of \(\Pi_t\) provides
these evaluation maps.  For \(u:c'\to c\), define its
restriction by \((b\cdot u)_a=b_{u a}\).  This is a map
of families: composition sends the incoming-arrow family at
\(c'\) into that at \(c\), and transposition along its
dependent product constructs the restriction.  Identity and
composition laws follow from the category equations and the
uniqueness of transposition.  A family map \(k:U_CE\to B\)
transposes to
\[
 E\to R_CB,\qquad e\mapsto(k(e\cdot a))_{a:d\to c}.
\]
It is an action map by associativity.  Conversely evaluation
at \(i(c)\) recovers \(k(e)\).  If \(\theta:E\to R_CB\)
is an action map, then its \(a\) coordinate at \(e\) is
the identity coordinate at \(e\cdot a\), by action
naturality, so the two constructions recover \(\theta\).
The inverse identities and their naturality prove
\(U_C\dashv R_C\).

\paragraph{Finite limits, colimits, and monomorphisms.}
The underlying category \(\mathcal E/C_0\) is a topos
by the dependent-product theorem.  Pullback along \(s,t\)
preserves finite limits directly, or because composition
\(\Sigma_s,\Sigma_t\) is its left adjoint.  It also
preserves finite colimits, because it has the proved right
adjoint \(\Pi_s,\Pi_t\).

For a finite diagram of actions, take its underlying limit
in \(\mathcal E/C_0\).  The pulled-back limiting cones
and the diagram's action maps induce a unique map
\(t^*E\to s^*E\).  Its identity and associative laws
hold after the limiting projections and hence hold by the
universal property.  A cone of action maps has the unique
underlying mediator, and the same cancellation by projections
makes that mediator an action map.  Thus \(U_C\) creates
finite limits, including the terminal family.

For colimits do the dual construction, using preservation
of the colimiting cocones by the required pullbacks.  A map
from \(t^*E\) is induced by the action maps on the diagram;
its target is \(s^*E\).  Identity and composition are
equalities of maps from the relevant pulled-back colimit,
so equality on all its cocone legs proves them.  Pullback
along the composable-arrow projections also preserves these
colimits, as just proved for every base map.  This justifies
the associative-law check on that domain as well.  The
colimit mediator is equivariant by the same cocone argument.
Thus finite colimits, including initial objects and arbitrary
parallel-pair coequalizers, are created by \(U_C\).

If an underlying map is monic, cancellation makes the action
map monic.  Conversely if an action map \(m:E\to F\)
is monic, two family maps \(B\rightrightarrows U_CE\)
equalized by \(U_Cm\) transpose under \(L_C\dashv U_C\)
to maps \(L_CB\rightrightarrows E\) equalized by \(m\).
They are equal, hence their inverse transposes are equal.
This proves that the underlying map is monic.  An underlying
inverse of an action isomorphism is equivariant by canceling
the original inverse squares, so \(U_C\) also reflects
isomorphisms.

\paragraph{Exponentials and the two inverse identities.}
Let \(H=(U_CF)^{U_CG}\) be the exponential in the slice
\(\mathcal E/C_0\).  Inside the underlying family of
\(R_CH\), impose the compatibility equation of the body
for every incoming \(a\), composable \(b\), and element
\(g\) of the indicated fiber of \(G\).  Explicitly these
variables form a finite-pullback family over \(U_CR_CH\);
the two sides are maps to the indicated fiber of \(F\).
Take their equality subobject and its universal image under
the family projection, using the already proved \(\forall\).
This constructs an actual subfamily \(V\hookrightarrow U_CR_CH\).
Base change for \(\forall\) ensures that its condition
holds after every parameter map.

It is stable under the cofree action.  For \(u:c'\to c\),
the restricted assignments satisfy
\[
 \phi_{u(a b)}(g\cdot b)
   =\phi_{(u a)b}(g\cdot b)
   =\phi_{u a}(g)\cdot b,
\]
so the action restricts to \(V\); its laws are inherited
from the cofree action.  Evaluation at identity gives an
action map \(V\times G\to F\): for \(u:d\to c\),
compatibility with \(a=i(c)\), \(b=u\) states
\[
             \phi_u(g\cdot u)=\phi_{i(c)}(g)\cdot u,
\]
which is precisely evaluation naturality.

For an action map \(h:K\times G\to F\), define its
transpose at \(e\in K_c\) by
\[
                \phi_a(g)=h(e\cdot a,g).
\]
First transpose the fiberwise map to \(H\), then use
\(U_C\dashv R_C\) to make this an action map into
\(R_CH\).  Naturality of \(h\) gives
\[
 h(e\cdot(a b),g\cdot b)
                =h(e\cdot a,g)\cdot b,
\]
so it factors through \(V\).  Evaluation at the identity
recovers \(h\).  Conversely, for \(k:K\to V\), the
transpose of its evaluated map has \(a\) coordinate equal
to the identity coordinate of \(k(e\cdot a)\).
Action naturality of \(k\) makes this exactly the old
\(a\) coordinate of \(k(e)\).  Extensionality of the
dependent products and the slice exponentials equates the
two maps.  These inverse identities prove uniqueness and
the exponential adjunction, natural in \(K\) and \(F\).

\paragraph{The internal sieve classifier.}
Apply the cofree construction to the family
\(C_0\times T\to C_0\).  Its fibers are assignments
\(v_a\) to incoming arrows.  Impose the inequalities
\(v_a\leq v_{a b}\) using the internal Heyting order
subobject and \(\forall\) along the composable-arrow
family.  This is an actual subfamily \(T_C\); it is
action-stable because \(u(a b)=(u a)b\).
The constant-top assignment defines an action map from
the terminal presheaf, and its underlying mono makes it
monic.  A sieve assignment is this truth assignment iff
its identity coordinate is true: one implication is
immediate, and the other follows by applying its inequality
with \(a=i(c)\) to every incoming \(b\).

For a mono of actions \(A\hookrightarrow E\), the
underlying mono has its slice characteristic map
\(\chi_A:U_CE\to C_0\times T\).  The cofree adjunction
gives the action map with coordinates
\(\chi_A(e\cdot a)\).  Closure of \(A\) under actions
implies its sieve inequalities: membership at \(a\)
entails membership at \(a b\).  Thus this map factors
through \(T_C\).  Its identity coordinate is
\(\chi_A(e)\), so its truth pullback is exactly \(A\),
by the preceding identity-coordinate characterization and
the underlying classifier.  Finite-limit creation proves
that this is also a pullback in the action category.

For uniqueness, let \(\psi:E\to T_C\) have the same
truth pullback.  On the family of pairs \((e,a:d\to c)\),
its \(a\) coordinate is the identity coordinate of
\(\psi(e\cdot a)\), by action naturality.  That coordinate
is true exactly on the subobject where \(e\cdot a\)
factors through \(A\), by the truth pullback and the
identity characterization.  Uniqueness of the base
classifier equates this coordinate predicate with
\(\chi_A(e\cdot a)\), as maps to \(T\), not only as
global truth values.  Dependent-product extensionality
equates the cofree assignments, and the mono into the
cofree family gives \(\psi=\chi_A^C\).

We have now constructed finite limits, finite colimits,
exponentials, and a unique subobject classifier.  Power
objects and their evaluation membership follow from the
same exponential/classifier bijections proved for the
original package.  Every internal assignment and condition
was built from finite-pullback families, dependent products,
and universal equality predicates, all previously proved.
When the object of objects is initial, its slice and action
categories are degenerate by strictness of the initial
object; the same constructions and universal properties
still hold.  No inhabited-family or external completeness
assumption was used.
\end{proof}
\end{annexitem}

\begin{annexitem}{annex:p10.lawvere-sheaves.internal-kan}{thm:p10-lawvere-sheaves-internal-kan}{The actual internal left and right extension adjunctions}
\begin{proof}
All indexing expressions below are finite fiber products of
\(C_0,C_1,D_0,D_1\) and the given family of \(A\).
Universal assignment families use the constructed \(\Pi\),
and universal compatibility uses the constructed \(\forall\).
The preceding theorem gives the required finite coequalizers
of actions.  None is an assumed external small-category limit.

\paragraph{Restriction and the right extension.}
Restriction is the family pullback along \(F_0\), with
\(x\cdot u=x\cdot F_1(u)\); functor identities make
this an action and make restriction a functor.
Let \(K\to D_0\) be the family of
\((c,b:F(c)\to d)\).  Pull \(U_CA\) back to \(K\)
along its \(c\) projection and apply \(\Pi_{K\to D_0}\).
This constructs all assignments \(a_{c,b}\).
Over that assignment family, form the family of triples
\((u:c\to c',b:F(c')\to d)\) and equate
\(a_{c,bF(u)}\) with \(a_{c',b}\cdot u\), using
\(\forall\) of the equality subobject.  It gives exactly
the subfamily in the body.  Restriction along
\(v:d'\to d\) is precomposition, \(b\mapsto vb\);
it respects the equation since \(v(bF(u))=(vb)F(u)\).
Transposition constructs this as a family map, and
associativity and identities prove its presheaf laws.

For an action map \(h:F^*B\to A\), define
\[
 \widehat h_d(x)_{c,b}=h_c(x\cdot b).
\]
This is a map into the assignment family by dependent
transposition.  For \(u:c\to c'\) and \(b:F(c')\to d\),
its compatibility follows from
\[
 h_c(x\cdot(bF(u)))
   =h_c((x\cdot b)\cdot F(u))
   =h_{c'}(x\cdot b)\cdot u.
\]
Restriction in \(D\) commutes with it because
\((x\cdot v)\cdot b=x\cdot(vb)\).
Thus \(\widehat h:B\to F_*A\) is an action map.

Conversely an action map \(\theta:B\to F_*A\) gives
\[
              h_c(x)=\theta_{F(c)}(x)_{c,i(F(c))}.
\]
For \(u:c\to c'\), action naturality of \(\theta\)
and its assignment equation give
\[
\begin{aligned}
 h_c(x\cdot F(u))
  &=\theta_{F(c')}(x)_{c,F(u)}\\
  &=\theta_{F(c')}(x)_{c',i(F(c'))}\cdot u
    =h_{c'}(x)\cdot u.
\end{aligned}
\]
So \(h:F^*B\to A\) is an action map.  Starting from
\(h\), identity evaluation gives \(h\) back.  Starting
from \(\theta\), action naturality along
\(b:F(c)\to d\) gives
\[
 \theta_{F(c)}(x\cdot b)_{c,i(F(c))}
                              =\theta_d(x)_{c,b},
\]
so transposition gives \(\theta\) back at every coordinate.
Dependent-product extensionality proves equality as maps.
These inverse identities are natural in \(A,B\), proving
\(F^*\dashv F_*\) and its unit/counit triangles.
On a map \(A\to A'\), postcompose each coordinate;
this preserves compatibility and gives the right-adjoint
functor's identity and composition laws.

\paragraph{The left extension is a finite coequalizer.}
The objects \(B_0,B_1\) have the actual finite-pullback
definitions displayed in the body.  Precomposition of
\(b:d\to F(c)\) gives their \(D\)-action; the action
laws are the category laws in \(D\).
Both parallel maps are action maps, since precomposition
commutes with \(F(u)b\).  Their coequalizer therefore
exists by the proved finite-colimit construction in
\(D^{\mathrm{op}}\mathcal E\).  Denote it \(F_0A\).

For a \(D\)-presheaf \(B\), maps \(F_0A\to B\)
are exactly action maps \(k:B_0\to B\) equalizing the
two parallel maps.  An action map \(k\) is determined
by its identity values
\[
                   h_c(x)=k(c,x,i(F(c))),
\]
since its naturality along \(b:d\to F(c)\) forces
\[
                        k(c,x,b)=h_c(x)\cdot b.
\]
These are actual family maps by the identity section of
\(D_1\) and the action map, so the claim is valid at
all internal parameter objects.
The parallel-pair condition becomes
\[
      h_c(x'\cdot u)\cdot b
                  =h_{c'}(x')\cdot(F(u)b).
\]
Taking \(b=i(F(c))\) gives exactly the naturality
condition for \(h:A\to F^*B\).  Conversely that
condition, followed by action along \(b\), gives
the displayed equation for every \(b\).
Thus the assignments
\[
        h\longmapsto k(c,x,b)=h_c(x)\cdot b,
        \qquad k\longmapsto h_c(x)=k(c,x,i(F(c)))
\]
are inverse and natural, and the coequalizer universal
property gives the required left-adjoint bijection.
Its unit is the coequalizer class of
\((c,x,i(F(c)))\); its counit evaluates a generator
\((c,x,b)\) as \(x\cdot b\).
The two proved inverse assignments give both triangle
identities.  A map of \(C\)-presheaves maps each
generator \((c,x,b)\) by its \(x\) coordinate and
respects the relations, so induces \(F_0\) on maps;
uniqueness of coequalizer mediators proves its identity
and composition laws.

Empty incoming-arrow families give terminal dependent
products in the right construction; empty generator
families give initial objects in the left construction.
All proofs use universal maps and equality predicates,
not the selection of global representatives of classes
or an assumption that epimorphisms split.
\end{proof}
\end{annexitem}

\begin{annexitem}{annex:p10.lawvere-sheaves.presheaf}{thm:p10-lawvere-sheaves-presheaf-bridge}{The exact concrete source-to-host bridge}
\begin{proof}
Use precisely the small category and fixed universe of
Definition~\ref{def:p10-presheaf-category}.  Theorem
\ref{thm:p10-presheaf-limits} constructs all small limits and colimits,
with restrictions induced by their universal properties.  In
particular it supplies both finite assumptions of the original
package, including coproducts and coequalizers.  No inference from
cartesian closedness alone is needed for finite colimits.

Theorem~\ref{thm:p10-presheaf-exponentials} gives
\(\operatorname{Nat}(\mathrm y c\times X,Y)\), the stated
restriction and evaluation, and the natural transpose
\[
 [\Lambda h_c(z)]_d(u,x)=h_d(Z(u)z,x).
\]
The inverse evaluation/transposition identities in that theorem are
exactly the adjunction required by the source's \((-)^X\); its
proof includes naturality at every arrow, not merely a componentwise
function-set formula.

Definition~\ref{def:p10-presheaf-sieves} and
Theorem~\ref{thm:p10-presheaf-topos} construct the sieve presheaf.
Renaming that already constructed object \(T=\Omega\) identifies
it with the crible object on p.~330.  To check the identification
explicitly, \(u\in(\chi_A)_c(x)\) means
\(X(u)x\in A(d)\).  Hence, for \(v:e\to d\),
\[
 v\in T(u)((\chi_A)_c(x))
 \ \Longleftrightarrow\ X(uv)x\in A(e)
 \ \Longleftrightarrow\ v\in(\chi_A)_d(X(u)x).
\]
This is naturality with the original variance.  The characteristic
sieve is maximal exactly when it contains \(1_c\), equivalently
when \(x\in A(c)\); thus its truth pullback is \(A\).
If another characteristic map \(\chi\) has that pullback, then
\[
 u\in\chi_c(x)
 \ \Longleftrightarrow\ 1_d\in\chi_d(X(u)x)
 \ \Longleftrightarrow\ X(u)x\in A(d),
\]
so it is the displayed map.  For a mono represented by a different
domain, Theorem~\ref{thm:p10-presheaf-mono-epi} gives its unique
image subpresheaf up to isomorphism over \(X\).

Finally Theorem~\ref{thm:p10-presheaf-power-objects} identifies
\(T^X(c)\) with \(\operatorname{Sub}(\mathrm y c\times X)\).
The restriction pulls back along \(\mathrm y u\times1_X\),
and membership is evaluation's truth pullback; at \(c\) it is
the condition \((1_c,x)\in S(c)\).  The bijections are natural
and classify each subobject uniquely.  All the reused results include
the empty base.  Their hypotheses and actual formulas therefore
supply the complete concrete package, while their detailed later
source attributions remain those of the host proofs.
\end{proof}
\end{annexitem}

\begin{annexitem}{annex:p10.lawvere-sheaves.choice}{thm:p10-lawvere-sheaves-choice}{Support, dependent product, and split epimorphisms}
\begin{proof}
We use image factorizations and the stated slice right adjoints.
First note that, for a subobject \(i:S\hookrightarrow X\),
\(\Pi_f(i)\) is a subobject of \(Y\).  Indeed \(i\) is a
monomorphism to the terminal object of \(\mathcal E/X\), and
the right adjoint preserves monomorphisms and terminal objects.
For a subobject \(B\hookrightarrow Y\), the slice adjunction
gives
\[
 B\leq\Pi_f(S)\ \Longleftrightarrow\ f^*B\leq S.
\]
By the predicate right adjunction, this is also
\(B\leq\forall_f S\).  Taking each of the two subobjects for
\(B\) proves \(\Pi_f(S)=\forall_f S\), including their
classified predicates.

Suppose every epimorphism splits.  Condition \((\exists)\)
is immediate for the epimorphism onto each object's support.
For any family \(p:E\to X\), write its image factorization as
\[
                 E\xrightarrow{q}S\xrightarrow{i}X.
\]
An underlying section \(r:S\to E\) of \(q\) is also over
\(X\), since \(pr=iqr=i\).  Applying \(\Pi_f\) gives a
section of \(\Pi_fq\).  Thus \(\Pi_fE\to\Pi_fS\) is a
split epimorphism in the slice, and its underlying map is a split
epimorphism.  The composite with the mono
\(\Pi_fS\hookrightarrow Y\) is an image factorization, so
\[
 \sigma_Y(\Pi_fp)=\Pi_fS=\forall_fS
                         =\forall_f\sigma_X(p).
\]
This proves \((\Pi)\).

Conversely suppose \((\Pi)\) and \((\exists)\) hold, and
let \(p:E\twoheadrightarrow X\) be any epimorphism.
Its support as an \(X\)-family is top.  Set \(f:X\to1\).
The right predicate adjoint preserves top, whence
\[
        \sigma_1(\Pi_fp)=\forall_f\mathsf{true}
                          =\mathsf{true}.
\]
Thus the object \(\Pi_fp\) has support \(1\).  Condition
\((\exists)\) gives \(1\to\Pi_fp\).  Transposition along
\(f^*\dashv\Pi_f\) gives a map from the terminal family
\(1_X:X\to X\) to \(p:E\to X\), namely a section of
\(p\).  Since \(p\) was arbitrary, every epimorphism splits.
No step turns an epi onto a support into a section over an arbitrary
base without using \((\Pi)\).
\end{proof}

The examples in Definition~\ref{def:p10-lawvere-sheaves-choice-examples}
separate the conditions concretely.  In \(G\mathsf{Set}\), for
\(f:X\to Y\) and \(p:E\to X\), the underlying fiber of
\(\Pi_fp\) over \(y\) is
\[
             \prod_{x\in f^{-1}(y)}p^{-1}(x).
\]
The action sends a section over \(y\) to the transported section
over \(gy\).  In external sets with choice this fiber is nonempty
exactly when every factor is nonempty, which is \((\Pi)\).
For an empty indexing fiber its product is a singleton, as universal
quantification requires.  The regular \(G\)-set is nonempty and
has no fixed point when \(G\ne1\), so its epimorphism to the
terminal \(G\)-set does not split; this refutes \((\exists)\).

For a covariant diagram \(E:P\to\mathsf{Set}\) with \(P\)
well-ordered, its support consists of stages where \(E(p)\) is
nonempty.  If this support is empty, the split map has empty domain
and codomain.  Otherwise it has a least stage \(p_0\).  Choose
one \(e_0\in E(p_0)\) and put
\(e_p=E(p_0\leq p)e_0\) for all supported \(p\).
The functor composition law proves compatibility of these elements;
they give a section over the support.  This uses one element, not
choice of an unrelated element at each stage.

For the specified arrow diagram, a section of \(p:E\to X\)
has forced values \(s_0(a)=e_a\), \(s_0(b)=e_b\).
Naturality would give both \(s_1(c)=u\) and \(s_1(c)=v\),
contradicting \(u\ne v\).  More explicitly, if \(R_0\) and
\(R_1\) are the covariant representables of the two stages, then
\(R_0=1\) and \(R_1=(\varnothing\to\{*\})\).
Adjunction and Yoneda identify the component of \(\Pi_{X\to1}p\)
at \(i\) with maps \(R_i\times X\to E\) over \(X\).
At \(0\) these are the nonexistent natural sections; at \(1\)
they are the two choices of \(u\) or \(v\).  Its support is
\((\varnothing\to\{*\})\), although \(\sigma_X(p)\)
is top and its universal image is top.  This verifies the exact
failure of the source square, not just failure of a section.
\end{annexitem}

\begin{annexitem}{annex:p10.lawvere-sheaves.closure}{thm:p10-lawvere-sheaves-closure}{Closure from the original modal-operator axioms}
\begin{proof}
Let \(m:A\hookrightarrow X\) have classifier \(\chi_A\).
Since \(\chi_Am=\mathsf{true}\,!_A\), the identity
\(j\mathsf{true}=\mathsf{true}\) gives
\(j\chi_Am=\mathsf{true}\,!_A\).  The classifier pullback
property therefore factors \(m\) through the subobject classified
by \(j\chi_A\).  This is extensivity; it is not an extra
unwritten axiom.

Intersections are classified by the conjunction of their two
predicates.  The meet identity gives
\[
        \overline{A\cap B}=\overline A\cap\overline B.
\]
If \(A\leq B\), then \(A=A\cap B\), so the displayed
identity implies \(\overline A\leq\overline B\).
Idempotence of \(j\) gives
\(\overline{\overline A}=\overline A\).
For \(g:Y\to X\), the classifier of \(g^*A\) is
\(\chi_Ag\), and
\(j(\chi_Ag)=(j\chi_A)g\).  Hence
\(\overline{g^*A}=g^*\overline A\).
The same proof applies to generalized predicates, since the
classifier identities are morphism identities.  Thus closure is
universal and preserves the terminal subobject as well as binary
meets.  No infinite colimit or external powerset of all objects
of \(\mathcal E\) was used.
\end{proof}
\end{annexitem}

\begin{annexitem}{annex:p10.lawvere-sheaves.associated-sheaf}{thm:p10-lawvere-sheaves-associated-sheaf}{Closed predicates, separated quotients, and the reflector}
\begin{proof}
All closures in this proof use
Theorem~\ref{thm:p10-lawvere-sheaves-closure}; all images and
effective quotients use their explicit constructions in
Theorems~\ref{thm:p10-lawvere-sheaves-predicate-adjoints}
and~\ref{thm:p10-lawvere-sheaves-effective-quotients}.
None is an extra hypothesis of sheafification.

\paragraph{Closed predicates and the first observation.}
Let \(i:T_j\hookrightarrow T\) be the equalizer of
\(j\) and \(1_T\).  Idempotence factors \(j\) as
\(i h\), and \(h i=1_{T_j}\) follows by canceling
\(i\) from \(i h i=j i=i\).  Thus \(h\) is a split
epi and \(i\) is the image of \(j\), by the proved
image uniqueness.  A map \(Z\to T\) factors through
\(T_j\) exactly when it classifies a closed subobject
of \(Z\).

Take a dense mono \(m:A\hookrightarrow B\) and a map
\(u:A\to T_j\), classifying a closed
\(U\hookrightarrow A\).  Regard \(U\) as a subobject
of \(B\) through \(m\), and take its closure in \(B\).
Pullback stability of closure gives
\[
              m^*(\overline U^{B})=\overline U^{A}=U.
\]
The closed-subobject classifier of \(\overline U^{B}\)
is therefore an extension \(B\to T_j\) of \(u\).
If two closed subobjects \(V,W\subseteq B\) agree
on \(A\), then
\[
 \overline{V\cap A}^{B}
      =\overline V^{B}\cap\overline A^{B}=V,
 \qquad
 \overline{W\cap A}^{B}=W.
\]
Their equality on \(A\) gives \(V=W\), and classifier
uniqueness gives uniqueness of the extension.  This proves
that \(T_j\) is a sheaf.  The argument uses arbitrary
subobjects and maps, so also applies after every generalized
parameter map and dense-mono pullback.

\paragraph{Exponentials and the second observation.}
If \(Y\) is a sheaf, transpose a map \(A\to Y^X\)
to \(A\times X\to Y\).  For any dense
\(A\hookrightarrow B\), its product with \(X\)
is dense, being its pullback along \(B\times X\to B\).
Extend uniquely to \(B\times X\to Y\), then transpose
back.  The two inverse transpose identities prove extension
and uniqueness for \(Y^X\).  Thus \(Y^X\) is a sheaf
without any inhabited-object assumption.

\paragraph{Separated objects and closed subobjects.}
An object \(Z\) has a closed diagonal iff maps into
\(Z\) have at most one extension across every dense
mono.  For one direction, the equalizer of two maps
\(B\rightrightarrows Z\) is closed, as a pullback
of the closed diagonal.  If it contains dense \(A\),
it is all of \(B\).  Conversely put
\(R=\overline{\Delta_Z}\subseteq Z\times Z\).
The mono \(\Delta_Z\hookrightarrow R\) is dense by
pullback stability of closure.  Its two coordinate maps
to \(Z\) agree on the diagonal, hence agree on \(R\)
by the uniqueness property.  This makes \(R\) factor
through the diagonal, so the diagonal is closed.
In particular every sheaf is separated, and every subobject
of a separated object is separated, since its diagonal is
a pullback of the ambient diagonal.

A closed subobject \(A\hookrightarrow Y\) of a sheaf
is a sheaf.  Extend a map from dense \(D\hookrightarrow B\)
to \(A\) first as a map \(B\to Y\).  Its inverse
image of \(A\) is closed and contains \(D\), hence
is all of \(B\).  The extension factors through \(A\);
its uniqueness follows from that in \(Y\).
If a sheaf \(A\) embeds in a separated object \(Z\),
let \(C=\overline A^{Z}\).  The dense mono
\(A\hookrightarrow C\) allows the identity of \(A\)
to extend to \(r:C\to A\).  The maps \(mr\) and
\(1_C\) agree on \(A\), and \(C\) is separated
as a subobject of \(Z\); uniqueness makes them equal.
Together with \(rm=1_A\), this proves \(A=C\).
Thus a sheaf is closed in every separated superobject.

\paragraph{Closed equality and the third observation.}
The relation \(R_X\) has classifier \(j\delta_X\).
It is reflexive by extensivity of closure and symmetric
by pullback along interchange of the two coordinates.
On \(X^3\), equality has
\(\delta_{12}\land\delta_{23}\leq\delta_{13}\).
Monotonicity and meet preservation give
\[
 j\delta_{12}\land j\delta_{23}
     =j(\delta_{12}\land\delta_{23})\leq j\delta_{13},
\]
which is transitivity.  These are inclusions of actual
subobjects at arbitrary parameter bases.  Thus \(R_X\)
is an internal equivalence relation, not just a relation
on global points.

Because \(j^2=j\), closed equality factors through
\(T_j\), and transposes to \(\kappa_X:X\to T_j^X\).
Equality of its two rows is precisely \(R_X\), by
the row argument already proved for effective quotients:
equivalent points have equal rows by symmetry/transitivity,
and equality of rows evaluated at the first point gives
equivalence by reflexivity and symmetry.  The inclusion
\(T_j\to T\) and its exponential are monic, so
computing row equality in \(T_j^X\) or \(T^X\)
gives the same kernel pair.  The image factorization
\(X\twoheadrightarrow S_X\hookrightarrow T_j^X\)
therefore gives an effective quotient of \(R_X\).

The ambient \(T_j^X\) is a sheaf by the first two
observations; its subobject \(S_X\) is separated.
For any \(f:X\to Z\) with \(Z\) separated, its
kernel relation is closed and contains \(\Delta_X\),
so contains \(R_X\).  The effective-quotient universal
property factors \(f\) uniquely through \(q_X\).
This proves the separated reflection and constructs the
quotient rather than assuming it as a preliminary step.

\paragraph{The associated sheaf and fourth observation.}
The object \(a_jX=\overline{S_X}^{T_j^X}\) is
a closed subobject of a sheaf, hence is a sheaf.
The mono \(S_X\hookrightarrow a_jX\) is dense by
pullback stability of closure.  A map \(X\to Y\)
to a sheaf first factors uniquely through \(q_X\),
since \(Y\) is separated, and then extends uniquely
from \(S_X\) to \(a_jX\), since \(Y\) is a
sheaf.  This gives the stated inverse bijections;
their uniqueness gives naturality in both variables.
On a map \(f:X\to X'\), transpose \(\eta_{X'}f\)
to define \(a_jf\).  The same uniqueness proves
identity/composition laws and naturality of \(\eta\).
For a sheaf \(Y\), the unit is an isomorphism by
this universal property, giving the full reflector
adjunction and its triangles.

If \(X\hookrightarrow Y\) is a mono with \(Y\)
a sheaf, its closure in \(Y\) is a closed subobject
of a sheaf and is therefore a sheaf.  The embedding
of \(X\) in its closure is dense.  Extension into
any sheaf is consequently unique and unrestricted,
which gives exactly the associated-sheaf universal
property.  This proves the fourth observation for
every such embedding, not just the chosen singleton
embedding used in the construction.

Finally, suppose \(Z\) is separated and closed in
every separated superobject.  Its closed-singleton
map \(\kappa_Z\) is monic, since its kernel pair
\(R_Z\) is now the diagonal.  Its target \(T_j^Z\)
is a sheaf, hence separated.  The hypothesis makes
this mono closed, so \(Z\) is a closed subobject
of a sheaf and therefore a sheaf.  Together with
the forward argument above this proves the exact
separated/closed characterization.  All steps used
the constructed internal maps, universal closure
and effective quotients; no infinite colimit,
Boolean negation or chosen class representative
was required.
\end{proof}
\end{annexitem}

\begin{annexitem}{annex:p10.lawvere-sheaves.sheaf-topos}{thm:p10-lawvere-sheaves-sheaf-topos}{Inverted maps, arbitrary pullbacks, and the sheaf classifier}
\begin{proof}
\paragraph{Precisely which maps are inverted.}
For \(f:X\to Y\), density of its relative diagonal
is equivalent to
\[
                    X\times_YX\leq R_X
                         \quad\text{in }\operatorname{Sub}(X\times X),
\]
by pulling the closure of \(\Delta_X\) back to that
kernel pair.  Suppose this and image density hold.
Factor \(f\) as \(X\twoheadrightarrow I\hookrightarrow Y\)
by the already proved stable image construction.  A map
\(h:X\to Z\) to a sheaf has a closed kernel relation,
which contains \(R_X\), so it coequalizes the kernel
pair of \(f\), equivalently of its epi onto \(I\).
The proved regular-epi universal property gives a unique
map \(I\to Z\).  Its dense mono into \(Y\) gives a
unique extension \(Y\to Z\).  Conversely two maps
on \(Y\) agreeing after \(f\) agree on \(I\) by
epi cancellation and then on \(Y\) by separatedness.
Thus \(f\) gives the bijection into every sheaf.

Suppose conversely that these bijections hold.  Let
\(J=\overline{\operatorname{Im}f}^{Y}\), with its
closed predicate \(c:Y\to T_j\).  Since \(f\)
factors through \(J\), its composites with \(c\)
and the constant truth map agree.  Injectivity of
the bijection for \(T_j\) equates those two maps
on \(Y\), so \(J=Y\): the image is dense.
Surjectivity for the sheaf \(T_j^X\) factors
\(\kappa_X\) through \(f\).  Therefore the kernel
pair of \(f\) lies in the kernel pair of
\(\kappa_X\), namely \(R_X\).  This gives density
of the relative diagonal.  By the reflector
adjunction, the bijections into every sheaf are
equivalent to \(a_jf\) being an isomorphism:
they are the natural hom bijections induced by
\(a_jf\), and taking its two source/target objects
gives an inverse by the usual hom universal property.
This proves all three equivalent conditions.

\paragraph{Stability and the units.}
Images commute with pullback by their proved construction,
and closure commutes with pullback by its proved law;
thus dense image is stable.  The relative diagonal
of a pulled-back map is the pullback of its original
relative diagonal, so its density is stable as well.
The unit \(\eta_X\) has image \(S_X\) dense in
\(a_jX\), and kernel pair \(R_X\) because its
second factor is monic.  The diagonal is dense
in \(R_X=\overline{\Delta_X}\).  Hence every
unit satisfies both conditions.

We shall also use the following image-density calculation.
If \(u:A\to B\) has dense image and \(v:B\to C\)
is arbitrary, then
\[
        \overline{\operatorname{Im}(vu)}^{C}
                  =\overline{\operatorname{Im}v}^{C}.
\]
Indeed the left side is closed and its inverse image
under \(v\) contains \(\operatorname{Im}u\), hence
is all of \(B\).  Image minimality puts
\(\operatorname{Im}v\) in the left side, while
the reverse inclusion follows from monotonicity.
This proves composition of dense-image maps, rather
than assuming a localization's closure properties.

\paragraph{The reflector preserves finite limits.}
An existing finite limit of sheaves in \(\mathcal E\)
is a sheaf: extend each component of a map across a
dense mono; the extensions are compatible because
they agree on that dense mono and their targets
are separated.  The base limit universal property
then gives the unique extension to the limit.
In particular \(1\) is a sheaf, so \(a_j1\cong1\).

For arbitrary maps \(f:X\to Z\), \(g:Y\to Z\),
put \(P=X\times_ZY\) and
\(Q=a_jX\times_{a_jZ}a_jY\).  The latter is a
sheaf by the preceding paragraph.  The product
\(\eta_X\times\eta_Y\) has dense image: its
image is \(S_X\times S_Y\), since products of
the two image epis are epi by the proved pullback
stability, and closure of that product is the
intersection of two pulled-back dense closures.

Let \(V\) be its inverse image of \(Q\) in
\(X\times Y\).  This is the relation
\(\eta_Zf(x)=\eta_Zg(y)\), by unit naturality.
The kernel pair of \(\eta_Z\) is \(R_Z\), so
universal closure gives
\[
                            V=\overline P^{X\times Y}.
\]
Thus \(P\hookrightarrow V\) is dense, and
\(V\to Q\) has dense image, being the pullback
of \(\eta_X\times\eta_Y\).  The image-density
calculation proves that the canonical \(k:P\to Q\)
has dense image.

Its kernel pair is \(\overline{\Delta_P}\).
Indeed on \(P\times P\) it says simultaneously
\(j\delta_X(x,x')\) and \(j\delta_Y(y,y')\).
Equality of pairs in \(P\subseteq X\times Y\)
is classified by \(\delta_X\land\delta_Y\);
meet preservation gives
\[
                j(\delta_X\land\delta_Y)
                           =j\delta_X\land j\delta_Y.
\]
These are equality of the kernel-pair subobjects,
so the relative diagonal of \(k\) is dense.
The already proved inverted-map characterization
makes \(a_jk\) an isomorphism.  Since \(Q\)
is a sheaf, this gives the canonical isomorphism
\[
                    a_j(X\times_ZY)
                              \cong a_jX\times_{a_jZ}a_jY.
\]
The construction is exactly the comparison induced
by the two reflected projections, because its
precomposition with \(\eta_P\) is \(k\).
Together with terminal preservation it proves
left exactness for all finite diagrams.  No
inhabitance, pointwise Boolean truth, or assumed
reflectivity entered the argument.

\paragraph{All structures of the sheaf topos.}
The full sheaf category has finite limits computed
in \(\mathcal E\), as proved above.  A finite
colimit is \(a_j\) applied to the base colimit:
maps from that reflected object to any sheaf are,
by adjunction and the base universal property,
exactly the compatible cocones of sheaf maps.
This includes the reflected initial object; it
need not equal the base initial object.
For sheaves \(X,Y\), the ambient \(Y^X\) is
a sheaf by the second observation.  Its ambient
evaluation and transpose bijections are therefore
the exponential bijections in the full sheaf
category, since its products are ambient products.

A mono between sheaves is also a base mono.  To
test cancellation of arbitrary maps \(W\to A\),
transpose them to maps \(a_jW\to A\) by the
reflector adjunction and cancel in the sheaf
category.  The converse follows by base
cancellation.  Such a mono is closed, since
its domain is a sheaf embedded in the separated
codomain.  Its base characteristic map factors
uniquely through \(T_j\), and its base truth
pullback is also the sheaf truth pullback by
finite-limit creation.  Conversely any sheaf
map \(X\to T_j\) classifies a closed base
subobject of \(X\), which is a sheaf because
it is closed in a sheaf.  Classifier uniqueness
in \(\mathcal E\), followed by monicity of
\(T_j\hookrightarrow T\), proves uniqueness
in \(\mathcal E_j\).  Thus
\(\mathsf{true}_j:1\to T_j\) is its full
subobject classifier, not merely a truth object.
Finite limits, finite colimits, exponentials,
and this classifier give the original topos
package.  The reflector is left exact and
left adjoint to the full inclusion, proving
the stated geometric morphism.

\paragraph{Every generalized parameter base.}
The proved slice package gives classifier \(B\times T\)
in \(\mathcal E/B\); its induced topology is
\(1_B\times j\).  Classification shows that its
closure of a mono is exactly the underlying base
closure.  Base and slice image monos coincide by
the already proved epi/mono criteria for slices,
and relative kernel pairs are the same underlying
fiber products.  Thus its dense-image and
dense-diagonal criteria are exactly the underlying
ones for maps over \(B\).

Let \(p:X\to B\) and put
\(W=B\times_{a_jB}a_jX\), with structure map to
\(B\).  This is a sheaf in the slice: a map
from a dense mono \(D\hookrightarrow C\) over
\(B\) consists of a map \(D\to a_jX\) whose
composite to \(a_jB\) equals \(\eta_B\) of
the parameter map.  Extend it uniquely to
\(C\to a_jX\).  The two composites to the
sheaf \(a_jB\) agree on dense \(D\), hence
agree on \(C\), so the extension factors
through \(W\) over \(B\).
The natural map \(X\to W\) is inverted by
\(a_j\): left exactness identifies \(a_jW\)
with \(a_jX\), and the induced map is the
canonical unit isomorphism there.  Its two
density conditions therefore hold, also in
the slice as just identified.  Apply the
already proved general inverted-map argument
in that slice to obtain unique extension into
every slice sheaf.  Since \(W\) is itself
a slice sheaf, it is the actual relative
associated sheaf of \(X\to B\).

For any \(b:B'\to B\), left exactness gives
\(a_j(B'\times_BX)\cong a_jB'\times_{a_jB}a_jX\).
The relative formula at \(B'\) consequently
becomes \(B'\times_{a_jB}a_jX\), which is
exactly the pullback of \(W\).  Their units
and extension maps agree by the same unique
universal property.  This proves canonical
base change at arbitrary generalized objects,
not just parameter maps with global sections.

For \(j=1_T\), dense monos are precisely
isomorphisms and the construction is the
identity reflection.  For \(j\) constantly
true, every mono is dense, including \(0\to B\);
the sheaf extension property makes every sheaf
terminal, and the sheaf topos is degenerate.
Both cases are included, illustrating why
no implicit nondegeneracy or unchanged initial
object may be inserted.
\end{proof}
\end{annexitem}

\begin{annexitem}{annex:p10.lawvere-sheaves.sheaves}{thm:p10-lawvere-sheaves-factorization}{The topology and actual geometric image factorization}
\begin{proof}
Write \(\mathcal B=\mathcal E''\),
\(\mathcal A=\mathcal E'\), and let
\(L=u^*:\mathcal B\to\mathcal A\) be left exact,
with right adjoint \(R=u_*\) and unit
\(\alpha_X:X\to RLX\).  Both categories have the
original topos package.  Their images, regular-epi
criteria and mono-epi isomorphism criteria therefore
come from the general constructions already proved,
not an extra geometric-factorization assumption.

\paragraph{The topology determined by the morphism.}
For a subobject \(A\hookrightarrow X\), let
\(L_XA=LA\hookrightarrow LX\).  For
\(B\hookrightarrow LX\), put
\[
               G_X(B)=X\times_{RLX}RB\hookrightarrow X,
\]
where the pullback uses \(\alpha_X\).  Right adjoints
preserve monos and limits, so this is a subobject.
The base adjunction gives
\[
                    L_XA\leq B\quad\Longleftrightarrow\quad
                                      A\leq G_X(B).
\]
Indeed a map \(LA\to B\) over \(LX\) transposes
to \(A\to RB\) whose composite into \(RLX\)
is \(\alpha_X\) composed with \(A\hookrightarrow X\);
this is exactly the pullback factorization.  The
inverse transpose proves the converse.  Thus this
is an actual poset adjunction on subobjects.

Define \(c_X=G_XL_X\).  The poset unit gives
extensivity; its counit gives
\(c_Xc_X\leq c_X\), and extensivity gives the
reverse inequality.  Both factors preserve finite
meets: \(L\) is left exact, and \(G_X\) is a
right adjoint.  They preserve the terminal
subobject as well.  For \(f:X\to Y\), unit
naturality and preservation of pullbacks by
\(L,R\) give
\[
          f^*G_Y(B)=G_X((Lf)^*B),
          \qquad f^*c_Y(A)=c_X(f^*A).
\]
Hence \(c\) is an idempotent meet-preserving
universal closure, constructed rather than assumed.

Let \(j_u:T''\to T''\) classify
\(c_{T''}(\mathsf{true})\).  For any subobject
with classifier \(\chi_A\), universal closure
gives \(c_XA=\chi_A^*c_{T''}(\mathsf{true})\),
so its classifier is \(j_u\chi_A\).
Preservation of top, meets, and idempotence of
\(c\), tested on the universal predicates,
therefore give
\[
       j_u\mathsf{true}=\mathsf{true},\qquad
       j_u\land=\land(j_u\times j_u),\qquad j_u^2=j_u.
\]
This constructs the source's single topology map.
For a mono \(m:A\hookrightarrow X\),
\[
             m\text{ is }j_u\text{-dense}
                 \quad\Longleftrightarrow\quad Lm\text{ is an isomorphism}.
\]
If \(L_XA\) is top then \(G_XL_XA\) is top.
Conversely if \(G_XL_XA\) is top, apply \(L_X\)
and its poset counit to get
\(\mathsf{top}\leq L_XA\), so \(L_XA\) is top.
This proves the equivalence without assuming that
\(L\) is conservative.

The functor \(L\) preserves image factorizations:
it preserves epis as a left adjoint and monos as
a left-exact functor, and image uniqueness was
proved in \(\mathcal A\).  It preserves kernel
pairs as well.  Consequently, for any map \(f\),
dense image means \(Lf\) is epic, and dense
relative diagonal means \(Lf\) is monic, by
the preceding criterion applied to those image
and diagonal monos.  Thus
\[
          Lf\text{ is an isomorphism}
                    \quad\Longleftrightarrow\quad
                           f\text{ is }j_u\text{-bijective}.
\]
This also proves the exact inverted-map criterion
needed for the factorization.

\paragraph{The two geometric factors.}
Every object \(RZ\) is a \(j_u\)-sheaf: for a
dense mono, extension maps into \(RZ\) correspond
by adjunction to maps in \(\mathcal A\) along
the isomorphism \(Lm\).  They extend uniquely.
Therefore \(R\) factors as
\[
           \mathcal A\xrightarrow{R'}\mathcal B_{j_u}
                                      \xrightarrow{i}\mathcal B.
\]
Let \(L'\) be the restriction of \(L\) to
\(\mathcal B_{j_u}\).  Full inclusion and the
original adjunction give \(L'\dashv R'\).
Finite limits of sheaves are base finite limits,
so \(L'\) is left exact.  It is the inverse
image of a geometric morphism
\(u':\mathcal A\to\mathcal B_{j_u}\).
The proved sheaf reflector \(a_{j_u}\dashv i\)
gives \(u'':\mathcal B_{j_u}\to\mathcal B\)
with full-and-faithful direct image \(i\).

Each sheafification unit \(\eta_X\) is
\(j_u\)-bijective by the preceding theorem,
so \(L\eta_X\) is an isomorphism.  These
natural isomorphisms identify \(L\) with
\(L'a_{j_u}\); the right-adjoint identification
is exactly \(R=iR'\).  Thus the adjoint pairs
give the stated factorization \(u=u''u'\).
If a map between \(j_u\)-sheaves is inverted
by \(L'\), it is \(j_u\)-bijective.  Its
sheafification is therefore an isomorphism,
and its units are already isomorphisms.  It
itself is an isomorphism.  This proves that
\(L'\) reflects isomorphisms, not merely that
the image functor has a chosen name.

\paragraph{Full-and-faithful morphisms and topology determination.}
If \(R\) was full and faithful, so is \(R'\),
and its counit is an isomorphism.  The triangle
identity makes \(L'\) of its unit an isomorphism;
the just-proved conservativity of \(L'\) makes
that unit an isomorphism too.  Hence
\(L'\dashv R'\) is an equivalence, identifying
the original domain with \(\mathcal B_{j_u}\)
over \(\mathcal B\).  It is the sheaf embedding
determined by the single map \(j_u\).

Conversely every topology has the constructed
sheaf embedding with left-exact reflector.
Its closure is recovered from its inverted
monos: if \(A\subseteq B\subseteq X\), then
\(A\hookrightarrow B\) is dense exactly when
\(B\subseteq\overline A^{X}\), by universal
closure.  Thus \(\overline A^{X}\) is the
largest intermediate subobject giving a dense
extension of \(A\).  The inverted monos of
the left adjoint determine these closures and
then determine \(j\) by the truth subobject.
This is the source's topology determination,
with categorical identifications up to the
canonical equivalences, not an unstated
uniqueness of strict chosen factorizations.
\end{proof}

The original p.~331 announces the factorization and gives
the finite closed-singleton route; the preceding proofs
complete those deductions editorially without upgrading
the historical source text's proof standing.  The following
source applications retain distinct proof obligations.

For the Godement example on pp.~331--332, a family of fibers over
\(X\) is sent to sections over the basis elements, while a
presheaf is realized by its local germs as a family over \(X\).
The source pairs these as \(\mathsf{lim}\dashv\Gamma\),
uses directness to make realization left exact, and obtains the
sheaf category as the geometric image.  Its alternative description
is a left-exact cotriple on \(\mathcal E/X\).  A full intrinsic
reconstruction is now supplied by
Theorem~\ref{thm:p10-lawvere-sheaves-godement}, deriving its
criterion as necessary and sufficient.  The qualified finite
Godement/filtering comparisons remain distinct from a formula
that the original article does not print.

The subsequent double-negation result on p.~332 is now proved
editorially in Theorem~\ref{thm:p10-lawvere-sheaves-double-negation}.
The original article announces it without the proof.  Classicality
refers to the sheaf category, not the original base or the stronger
source Skolem/choice package.
\end{annexitem}

\begin{annexitem}{annex:p10.lawvere-sheaves.godement}{thm:p10-lawvere-sheaves-godement}{Local directness, germs, sheaves, and the actual cotriple}
\begin{proof}
Write \(\mathcal P=A^{\mathrm{op}}\mathcal E\) and
\(\mathcal D=\mathcal E/X\).  Both topos packages,
dependent products, stable images and effective quotients
have already been constructed.  In particular no external
infinite colimit or set of global points is assumed.

\paragraph{The two functors and their actual adjunction.}
The internal order is the pullback of subobject inclusion
along \(A\to T^X\), so it is reflexive and transitive.
It is a thin category even when duplicate indices prevent
antisymmetry.  The first coequalizer object is
\(H_0(F)=P\times_AU_AF\) over \(X\).
The second is the pullback of the order-arrow object with
\(P\) at its source and \(U_AF\) at its target.
Order inclusion gives membership at that target.  Thus
both displayed generator maps are actual maps over \(X\).
Their finite coequalizer constructs \(L(F)\).  On an
action map, map its value coordinate in both objects;
the coequalizer universal property gives \(L\) on maps,
with identity and composition laws by uniqueness.

For \(M\to X\), dependent product along \(P\to A\)
constructs \(\Gamma M\).  An order arrow \(a\leq b\)
gives the inclusion \(U_a\hookrightarrow U_b\);
precomposing a section over \(U_b\) gives a section over
\(U_a\).  Dependent transposition and its uniqueness
construct this action map and prove its identity and
composition laws at every parameter object.
Maps \(L(F)\to M\) correspond to maps
\(H_0(F)\to M\) equalizing the two generators.
Dependent transposition identifies the latter with
family maps \(F\to\Gamma M\); the parallel-pair
equation is exactly their action compatibility.
These correspondences are inverse and natural by the
coequalizer and dependent-product universal properties.
The unit sends \(v\in F_a\) to its section
\(x\mapsto[x,a,v]\); the counit evaluates a germ
of a section at \(x\).  Both inverse identities
give the adjunction triangles.  This proves the
adjunction independently of directness.

\paragraph{Equality of germs under local directness.}
On \(H_0(F)\times_XH_0(F)\), take the image relation
\(R_F\) of the witness object
\[
 (x,a,b,c,s,t):\quad P(x,c),\ c\leq a,b,
                 s\in F_a,\ t\in F_b,\ s|c=t|c.
\]
It is reflexive by \(c=a\), and symmetric by
interchange of the two values.  For transitivity,
pull the witness epis back to a pair of composable
relations.  They give witnesses \(c,d\) locally.
The common-refinement epi in the directness condition,
pulled back to these data, gives locally \(e\leq c,d\)
containing \(x\).  Action composition restricts both
equalities to \(e\), proving the composite relation.
The composite of these actual witness epis is epic,
so image inclusion proves transitivity on the whole
parameter object.  This is an internal equivalence
relation, not a choice of one witness per global point.

It contains the generator pairs: use their source index
as the common refinement.  Conversely a witness at
\(c\) gives two generator identifications with the
same restricted value at \(c\), so the coequalizer
map is equal on that witness.  Cancel the witness epi
to conclude equality on \(R_F\).  The effective
quotient of \(R_F\) therefore coequalizes the original
parallel pair, and the original coequalizer coequalizes
\(R_F\).  Their universal properties give inverse
isomorphisms.  Its kernel pair is exactly \(R_F\).
This establishes the stated local germ-equality
criterion without a finite-path or transitive-closure
construction hidden in the notation.

\paragraph{Directness implies left exactness.}
For the terminal presheaf, representatives are just
members \(a\) containing \(x\).  The cover epi
\(P\to X\) gives local representatives.  Common
refinement relates any two, so its kernel relation is
\(P\times_XP\).  Since \(P\to X\) is a regular
epi by the proved exactness construction, its kernel
pair coequalizer is \(X\).  Thus \(L1_A\cong X\).

The comparison \(L(F\times G)\to LF\times_XLG\)
is epic.  Pull back the representative epis for the
two germs; locally they are \((a,s)\) and \((b,t)\)
over the same \(x\).  Pull back the refinement epi
to obtain \(c\leq a,b\) containing \(x\).  The
representative \((c,s|c,t|c)\) maps to the given
pair.  This gives an actual epic witness map into
the comparison's image, hence epimorphy.
For monicity, pull back representatives for two
product germs whose two projected germs agree.
The exact kernel descriptions for \(F\) and \(G\)
give local equality witnesses \(c,d\).  Refine them
locally to \(e\); both coordinates then agree at
\(e\), so the product germs agree.  These witness
maps are epis by stable image/pullback and directness.
Cancelling them equates the two projections of the
comparison's kernel pair, proving monicity.  The
mono-epi criterion makes the comparison an isomorphism.

Realization preserves monomorphisms: equality of
germs of an underlying subpresheaf in the ambient
presheaf is witnessed at a common refinement, where
component monicity cancels the two values.  The same
witness therefore gives equality within the subpresheaf.
Again representative/witness epis cancel to give
monicity of the realized map.
For parallel maps \(f,g:F\rightrightarrows G\), let
\(E\hookrightarrow F\) be their actual component
equalizer presheaf.  A germ in the equalizer of
\(Lf,Lg\) has a representative \((a,s)\) locally.
Equality of its two images has a common-refinement
witness \(c\leq a\) where
\(f(s)|c=g(s)|c\).  Naturality puts \(s|c\) in
\(E_c\), giving a germ in \(LE\) representing the
original germ.  The same epic-witness argument proves
the equalizer comparison epic; the preceding mono
argument makes it monic.  It is an isomorphism.
Terminal objects, binary products and equalizers
generate finite limits, proving left exactness.

\paragraph{Left exactness forces precisely these conditions.}
The coequalizer \(P\to L1_A\) is always epic.
If \(L\) preserves the terminal object, then
\(L1_A\cong X\), so \(P\to X\) is epic.  This
proves the cover condition rather than inserting it.

For an internal index \(a\), the representable
presheaf \(\mathrm ya\) has a value at \(c\)
exactly when \(c\leq a\).  Its realization is
\(U_a\): every representative \((x,c,c\leq a)\)
is identified by a generator with \((x,a,i(a))\),
and membership at \(c\) implies membership at
\(a\).  These maps give inverse isomorphisms,
without a directness assumption.  The image of
\(L(\mathrm ya\times\mathrm yb)\to X\) is
exactly the image of the family of common lower
members containing \(x\), by its representative epi.
If \(L\) preserves products, this is
\(U_a\times_XU_b\).  The common-refinement
projection is therefore epic.

This argument tests arbitrary generalized indices,
not just indices over \(1\).  To justify that step
explicitly, for a parameter object \(B\), the
constant presheaf \(\mathbf B\) has realization
\(L1_A\times B\), because product with \(B\)
preserves the defining finite coequalizer in the
slice.  Under terminal preservation this is
\(X\times B\).  The construction on presheaves
over \(\mathbf B\) is the corresponding slice
of \(L\); if \(L\) is left exact, that slice
functor is left exact.  Its representative and
coequalizer objects are exactly the original
objects pulled to that parameter base.  Apply
the product argument to the universal pair
\(a,b\) over \(A\times A\).  It gives the
single common-refinement epi in the definition.
Necessity and sufficiency are now both established,
so the criterion is not an arbitrary added premise.

\paragraph{The actual covering topology and image sheaves.}
For a sieve \(S\hookrightarrow\mathrm ya\), its
realization is the subobject of \(U_a\) given by
the image of \(\{(x,b):b\in S,\ P(x,b)\}\to X\).
This follows from its representative epi and
\(L\)'s now-proved preservation of monos.
The geometric-factorization theorem therefore says
that \(S\) is dense for the image topology
precisely when this subobject is all of \(U_a\),
which is the displayed coverage.
The maximal sieve covers because it contains \(a\).
If a sieve covers \(a\) and \(b\leq a\), its
pullback covers \(b\): for a point in \(U_b\),
the cover gives a member \(c\) locally, and
common refinement of \(b,c\) gives a member of
the pulled-back sieve.  Composing these epis
proves coverage everywhere.  For transitivity,
a covering sieve first gives a local member,
and its covering refinement gives a further
local member; composition of the witness epis
gives the composite covering sieve.  Thus these
are internal Grothendieck-covering axioms with
their actual local meaning, not global choices.

More explicitly, the induced topology closes a
sieve \(S\) on \(a\) by including \(b\leq a\)
exactly when \(U_b\) is contained in that image
union.  Indeed the subobject adjunction used
in geometric factorization gives
\(\mathrm yb\leq\overline S\) iff
\(L\mathrm yb\leq LS\).  The internal Yoneda
bijection here is just evaluation at identity,
with inverse given by action along each incoming
arrow; the same proof holds at every parameter
base.  It therefore identifies the topology
on the full internal sieve classifier.

Its sheaves are exactly presheaves with unique
extension from every covering sieve at every
generalized index.  One direction is the dense
mono extension property already proved for
the induced topology.  For the converse, let
\(M\hookrightarrow N\) be a dense mono of
presheaves and \(h:M\to F\) a map to a presheaf
having that covering-sieve property.  At an
element \(n\) of \(N_a\), its characteristic
sieve is the arrows on which its restriction
lies in \(M\); density makes it covering.
The map sending such a restriction to its
\(h\)-image extends uniquely from that sieve
to \(\mathrm ya\).  Take this construction
at the universal element of the underlying
family of \(N\), using that family itself
as parameter object.  Its extension is an
actual map of underlying families \(N\to F\),
not a selection of values for global elements.
Action naturality follows because restricting
the characteristic sieve gives the characteristic
sieve of the restricted element, and the two
extensions then agree by uniqueness.  Their
parameter naturality and uniqueness prove
the full dense-mono extension property.
This identifies the geometric image with the
usual internally defined basis-sheaf category.

\paragraph{The left-exact cotriple and its comparison.}
Factor the geometric adjunction as already proved,
with image topos \(\mathcal H\), conservative
left-exact \(L':\mathcal H\to\mathcal D\)
and right adjoint \(\Gamma'\).  The original
\(\Gamma\) lands in \(\mathcal H\), so
\(L\Gamma\cong L'\Gamma'\), including the
adjunction structure maps.  Let \(\zeta\)
be the unit of this restricted adjunction,
\(\varepsilon\) its counit, and
\(G=L'\Gamma'\).  It is left exact because
\(L'\) is left exact and \(\Gamma'\) preserves
limits.  Its comultiplication is
\(\delta_M=L'\zeta_{\Gamma'M}\); unit
naturality and the adjunction triangles prove
the comonad equations.  Thus this is the
source's actual left-exact cotriple on families.

For a coalgebra \(\theta:M\to GM\), with
\(\varepsilon_M\theta=1_M\) and
\(G\theta\,\theta=\delta_M\theta\), form
in \(\mathcal H\) the finite equalizer
\[
             N=\operatorname{Eq}
                    (\Gamma'\theta,\zeta_{\Gamma'M})
                            \hookrightarrow\Gamma'M.
\]
Applying \(L'\) gives the equalizer of
\(G\theta\) and \(\delta_M\).
The coalgebra map \(\theta\) itself is that
equalizer.  If \(r:Z\to GM\) equalizes
the pair, counit naturality and the triangle
give
\[
 \theta\varepsilon_Mr
     =\varepsilon_{GM}G\theta\,r
     =\varepsilon_{GM}\delta_Mr=r.
\]
Since \(\theta\) is split monic, this is
the unique factorization.  Thus \(L'N\cong M\).
This is a coalgebra isomorphism: write its equalizer
inclusion as \(i:N\to\Gamma'M\) and its isomorphism
as \(e=\varepsilon_M L'i\).  Unit naturality and
the right-adjoint triangle give
\(Ge\,L'\zeta_N=L'i=\theta e\).
It therefore respects the coalgebra structures, proving essential
surjectivity of the comparison by construction.

For \(S\in\mathcal H\), its unit \(\zeta_S\)
is the equalizer of
\(\Gamma'(L'\zeta_S)\) and
\(\zeta_{\Gamma'L'S}\).  It factors through
that equalizer by unit naturality; after \(L'\)
it is the preceding split coalgebra equalizer.
The induced map is therefore inverted by
\(L'\), and conservativity makes it an
isomorphism.  This proves the equalizer assertion
without assuming comparison full faithfulness.
If \(f:L'S\to L'T\) is a coalgebra morphism,
transpose it to \(\widehat f=\Gamma'f\,\zeta_S\).
The coalgebra equation and unit naturality give
\[
       \Gamma'(L'\zeta_T)\widehat f
                         =\zeta_{\Gamma'L'T}\widehat f.
\]
It factors uniquely through the just-proved
equalizer \(\zeta_T\), giving \(S\to T\).
Inverse transposition identifies its image
under \(L'\) with \(f\).  Conversely every
map \(S\to T\) respects the coalgebra
structures by unit naturality.  This proves
full faithfulness.  The equalizers and their
mediators are natural on coalgebra morphisms,
so they give an inverse comparison functor.
Together with essential surjectivity this
proves \(\mathcal H\simeq\operatorname{Coalg}(L\Gamma)\).
Only finite equalizers and proved conservativity
were used; no comonadicity theorem was a premise.

\paragraph{The ordinary Godement comparison.}
In sets with a topological basis, the coequalizer
is the germ set over \(X\).  Give it the
topology generated by the images of the germ
sections \(U_a\to LF\), one for each
\(v\in F_a\).  Equality on a common lower
member makes the intersection of two such
images a union of images on refinements.
Each image projects bijectively and homeomorphically
to \(U_a\), so the projection is a local
homeomorphism.  This uses exactly the cover
and refinement conditions, not an extra choice
of neighborhoods.  Continuous sections are
locally represented by the original sections.
For a basis sheaf, equality of germs implies
local equality and then equality by its
covering uniqueness axiom; compatible local
representatives glue by its existence axiom.
Thus its sections coincide with the original
basis sheaf, and an arbitrary presheaf gives
its generated sheaf of continuous sections.
These are the steps of Godement's inspected
II\S1.2, pp.~110--112; the intrinsic proof
above lifts them to all parameter objects
without importing setwise point sufficiency
into an arbitrary topos.
\end{proof}
\end{annexitem}

\begin{annexitem}{annex:p10.lawvere-sheaves.double-negation}{thm:p10-lawvere-sheaves-double-negation}{The internal Heyting proof and the canonical two-value classifier}
\begin{proof}
\paragraph{The topology identities are intuitionistic.}
Work in an arbitrary subobject Heyting algebra; all identities
are uniform in the parameter object and therefore are internal
classifier identities as well.  Residuation gives
\(p\land\neg p=0\) and \(p\leq\neg\neg p\).
Negation reverses order: if \(p\leq q\), then
\(\neg q\land p=0\), so \(\neg q\leq\neg p\).
Applying it to extensivity gives
\(\neg\neg\neg p\leq\neg p\), while extensivity
at \(\neg p\) gives the reverse inequality.  Thus
\(\neg^3p=\neg p\) and \(\neg^4p=\neg^2p\),
which proves idempotence.  Residuation also gives
\(\neg0=1\), \(\neg1=0\), so double negation
preserves both truth and bottom.

Monotonicity gives
\(\neg\neg(p\land q)\leq\neg\neg p\land\neg\neg q\).
For the reverse inequality put \(c=\neg(p\land q)\).
Then \(c\land p\leq\neg q\); hence
\(\neg\neg q\land c\land p=0\), so
\(\neg\neg q\land c\leq\neg p\).
Intersecting with \(\neg\neg p\) gives zero,
and residuation gives
\[
       \neg\neg p\land\neg\neg q\leq\neg c
                                      =\neg\neg(p\land q).
\]
This proves meet preservation without Boolean reasoning.
The previously constructed topology/sheaf machinery therefore
applies to \(j=\neg\neg\).

\paragraph{Closed predicates have actual complements.}
Triple negation makes \(\neg p\) fixed under double
negation.  For any two predicates,
\(\neg(p\lor q)=\neg p\land\neg q\): a predicate
\(r\) lies below the left side exactly when
\(r\land p=0=r\land q\), by distributivity and
residuation, which is the right-side universal property.
Consequently
\[
       \neg(p\lor\neg p)=\neg p\land\neg\neg p=0,
       \qquad \neg\neg(p\lor\neg p)=1.
\]

For a sheaf \(X\), its subobjects in the sheaf category
are precisely its closed ambient subobjects, by the full
classifier theorem.  Their least upper bound is the closure
of their ambient join: it is closed and contains them, and
every closed upper bound also contains that closure.
Their meet is the ambient meet.  Here bottom is ambient
zero, because \(j0=0\) and a closed subobject of a sheaf
is a sheaf.  Thus a closed \(A\) has closed complement
\(\neg A\), with meet zero and sheaf join top.
It is the Heyting negation in this closed-subobject lattice
because it is already the largest ambient predicate with
zero meet.  The displayed excluded-middle identity holds
for every sheaf parameter object and hence for its internal
classifier \(T_j\), not just for global truth values.
The relative-slice theorem supplies the same statement
after every generalized base change.

\paragraph{A disjoint cover is a genuine coproduct.}
We use only consequences of the original package already
proved for any topos, including the sheaf topos.  Suppose
\(A,B\hookrightarrow X\) are subobjects with
\(A\cap B=0\) and \(A\lor B=X\).
Let \(e:A+B\to X\) be their coproduct map.
Its image is their join, so it is epic.
Products in the slice over \(X\) preserve finite
colimits by its proved cartesian closedness.  Therefore
the kernel pair of \(e\) decomposes as
\[
 (A+B)\times_X(A+B)
 \cong(A\times_XA)+(A\times_XB)
        +(B\times_XA)+(B\times_XB)
 \cong A+0+0+B.
\]
Both kernel-pair projections agree on the \(A\) and
\(B\) parts and trivially on the initial parts.
They agree on the coproduct, so \(e\) is monic.
The proved mono-epi criterion makes it an isomorphism.
This derives a summand decomposition; no unproved
extensivity or disjoint-coproduct axiom is assumed.

\paragraph{All three classicality conditions are equivalent.}
If the classifier is internally Boolean, every subobject
has its Boolean complement.  The preceding argument gives
\(A+(\neg A)\cong X\).  Complements are unique:
if \(B,C\) both complement \(A\), distributivity gives
\(B=B\land(A\lor C)=B\land C\), and symmetrically
\(C=B\land C\).  Thus they are the same subobject,
with the canonical isomorphism between representatives.

Conversely if every mono has a complementary summand,
each subobject lattice is Boolean.  Indeed a complement
\(B\) of \(A\) is its Heyting negation: whenever
\(C\land A=0\), distributivity gives
\(C=C\land(A\lor B)=C\land B\), so \(C\leq B\).
Thus excluded middle holds in every parameter object;
testing universal predicates proves the internal Boolean
identities of the classifier.

For the classifier object itself, the truth subobject
\(1\hookrightarrow T\) is classified by \(1_T\).
The false subobject \(1\hookrightarrow T\) is
classified by its negation.  To verify the latter rather
than assume it, on any parameter object \(W\),
\(\neg p=\mathsf{true}\) iff \(p=\mathsf{false}\),
by residuation \(p\leq0\); this is exactly the
factorization criterion for the pullback of false.
Classifier uniqueness gives the stated predicate.
Their intersection is initial: on its pullback truth
equals false, so the identity subobject equals the
initial subobject there, forcing that object to be
initial by classification.  If the classifier is
Boolean these two subobjects join to all of \(T\).
The disjoint-cover result consequently makes their
canonical map \(1+1\to T\) an isomorphism.

Conversely, if that canonical map is an isomorphism,
pull its coproduct decomposition back along the
classifier \(X\to T\) of any mono.  Pullback
preserves finite colimits in slices by the proved
dependent-product adjoint.  Its two pulled-back
pieces are the given subobject and the false
piece, with initial intersection, so they form
the required complementary summands of \(X\).
This proves the remaining implication.
Apply these equivalences in the already constructed
double-negation sheaf topos to obtain precisely
\(1+1\cong T_j\).  This is classical logic only;
the separate support-splitting/Skolem assumption
does not occur in any proof step.
\end{proof}
\end{annexitem}

\begin{annexitem}{annex:p10.lawvere-sheaves.ultraproduct}{thm:p10-lawvere-sheaves-ultraproduct}{The exact localization assertion and ordinary-ultrafilter comparison}
\sketchnote{The original article announces the combined theorem without
complete proofs and without printing a precise closed-base convention.
The ordinary fraction calculus, localized topos and logical preservation
are proved in Annex items\ \ref{annex:p10.lawvere-localization.denominators},
\ref{annex:p10.lawvere-localization.fractions} and
\ref{annex:p10.lawvere-localization.logical}. The actual full-datum
partial-element closed functor and its hom/composition comparisons are
proved in \ref{annex:p10.lawvere-closed-base.partials} and
\ref{annex:p10.lawvere-closed-base.closed}; the countercases are
\ref{annex:p10.lawvere-closed-base.countercases}, and the nondegenerate
BVM collapse with its required homomorphism and choice is
\ref{annex:p10.lawvere-closed-base.collapse}. What is not supplied is an
arbitrary strong or normal identification over the original base, which
the full datum does not entail. The degenerate exception to nondegenerate
two-valued collapse is explicit; the historical sketch status does not
make the completed, qualified editorial constructions missing again.}
The datum is the specified finite-limit-preserving \(u_*\)
together with the specified Heyting homomorphism \(h\).  The
inversion class is not all epimorphisms and not all maps with
nonempty fibers; it is the monos with the exact universally
quantified truth value sent by \(h\) to top.  Localizing this
class is the construction named on p.~332. The independent ordinary
localization proof constructs the fractions and their equality/coherence,
then the finite limits/colimits, exponentials, classifier and quantifiers.
It does not infer these properties merely from the name of the construction,
and does not assert that the original base functor descends unchanged.
The new base functor instead quotients h-total partial members, preserving
the full internal \(u_* / h\) datum. Its ordinary fraction names inject
into the base hom objects. They need not exhaust their global elements,
as the diagonal countercase shows. In the particular BVM collapse over
the abstract-set base, the geometric original adjunction and base choice
do prove the additional normal element identification. This distinction
is part of the proof scope, not an assumed stronger convention.

The comparison \(\mathsf{Set}^I\to\mathsf{Set}\) can be
computed without that theorem.  Product preserves finite limits.
The truth-value object of \(\mathsf{Set}^I\) is the constant
family of two-element sets, so its product is \(2^I\).  A given
ultrafilter \(U\) gives a Boolean, hence Heyting, homomorphism
\[
       h(b)=1\quad\Longleftrightarrow\quad
                         \{i:b_i=1\}\in U.
\]
For a mono \((A_i\hookrightarrow X_i)_{i\in I}\), a section
of that mono at index \(i\) exists exactly when \(A_i=X_i\).
Its universally quantified support is therefore the truth vector
of this equality set, and the source inversion condition is
\(\{i:A_i=X_i\}\in U\).  This verifies the displayed
specialization.  It neither constructs an ultrafilter nor claims
that the category localization is an ordinary product quotient
of one chosen family of underlying sets.  The latter requires
its own partial-section and model-theoretic comparison argument.
\end{annexitem}

\begin{annexitem}{annex:p10.lawvere-sheaves.generators}{thm:p10-lawvere-sheaves-generators}{Why the source's Skolem and classical assumptions give generators}
\begin{proof}
Let \(f,g:E\to F\) be unequal.  Their equalizer
\(K\hookrightarrow E\) is not all of \(E\).  By the stated
complement hypothesis write \(E=K+D\); then \(D\ne0\).
Let \(D\twoheadrightarrow U\hookrightarrow1\) be its support
factorization.  If \(U=0\), there would be a map \(D\to0\).
In a topos the initial object is strict, so \(D=0\), a
contradiction.  Here strictness follows from cartesian closedness:
product with \(D\) preserves the initial object, and a map
\(D\to0\) makes \(D\times0\to D\) admit a section; a
retract of the initial object is initial.  Thus \(U\ne0\).

By \((\exists)\), the epi \(D\to U\) has a section
\(s:U\to D\).  Compose it with \(D\hookrightarrow E\)
to obtain \(a:U\to E\).  If \(fa=ga\), then \(a\)
factors through \(K\).  It already factors through \(D\),
so it factors through the disjoint intersection \(K\cap D=0\).
Strictness would imply \(U=0\), which was excluded.  Therefore
\(fa\ne ga\).  This proves that the subterminal objects
distinguish every unequal pair, the generating-family assertion.
The proof does not assume \(1\) itself is a generator or that
there are only two subobjects of \(1\).
\end{proof}
\end{annexitem}

\begin{annexitem}{annex:p10.lawvere-sheaves.ch}{thm:p10-lawvere-sheaves-ch}{The full extent and limits of the source's CH sketch}
\sketchnote{The original article leaves suitable undefined and announces
the no-epi and comparison conclusions without complete proofs. The
explicit sufficient internal ccc and actual finite-condition verification
are now proved in Annex items\
\ref{annex:p10.lawvere-ch-chain.cardinal-tools} and
\ref{annex:p10.lawvere-ch-chain.ccc}; the internal-site/constant-section
bridge, complete parametrized no-epi argument, generic injection descent
and both strict internal comparisons are
\ref{annex:p10.lawvere-ch-chain.constant-sections},
\ref{annex:p10.lawvere-ch-chain.no-epi},
\ref{annex:p10.lawvere-ch-chain.generic} and
\ref{annex:p10.lawvere-ch-chain.comparisons}.
These are independently supplied editorial proofs at the declared base
hypotheses, not an unprinted original definition of suitable. Specified
set-theoretic transfer/model interpretation and CH's positive consistency
direction remain unproved. The historical source sketch status does not
turn the six completed proof dependencies back into missing ones.}
The source's base \(\mathcal S\) includes choice, a nontrivial
recursion object, classicality, and irreducibility.  Its
\emph{Boolean} definition for \(\mathcal E\) includes the
Skolem condition as well as classicality; this must not be replaced
by the weaker modern use of ``Boolean topos.''  The claimed
equivalence to \(\mathcal S\)-complete Boolean algebras now has the
actual editorial functor, morphism, 2-cell and comparison proofs under
their expressly declared cell regimes. Neither those proofs nor the
further choice proof is printed on p.~333. Likewise the monoid
example distinguishes two-valuedness
of \(\operatorname{Sub}(1)\) from Boolean internal logic; it is
not a counterexample to the implication with both stated assumptions.

For the refutation construction, conditions are finite compatible
assignments on \(\omega\times I\), stronger by extension.
Every two distinct indices in \(I\) can have distinct binary
sequences witnessed by a further condition: choose a fresh natural
coordinate not already used for those indices and assign different
bits.  Conditions deciding one coordinate form a dense set.  These
are the local data motivating the source's injection
\(u^*I\hookrightarrow2^\omega\). The actual dense evaluation
subpresheaf, reflected inverse, transpose and fresh-coordinate
kernel argument are now \ref{annex:p10.lawvere-ch-chain.generic}.
That proof handles arbitrary base parameters; the article's description
of finite conditions alone was not its full descent proof.

The transported ground power object \(u^*(2^\omega)\) and the
new power object \(2^\omega\) have different roles.  A geometric
inverse image preserves finite limits, not necessarily this
exponential, so replacing the latter by the former would defeat the
construction.  The source's no-epi lemma is meant to prevent
collapse of the chosen ground cardinal comparison.  Applying it
retains the hypothesis \(Y\times\omega\cong Y\), the correct
ground objects and the expressly chosen sufficient internal ccc.
Its deciding-antichain proof and cardinal constructions are
\ref{annex:p10.lawvere-ch-chain.no-epi} and
\ref{annex:p10.lawvere-ch-chain.cardinal-tools}; the finite-input branch
is essential, since a ground epi from \(X\times\omega\) does not
in general give one from \(X\). Internal cofinal-witness/all-refinement
clauses are proved in \ref{annex:p10.lawvere-ch-chain.constant-sections},
not inferred by substituting an external set-valued site for this base.
The lemma's object \(\operatorname{Epi}(X,Y)\) concerns all
internally parametrized maps with full image, not just absence of
one selected global function.  The source's strict inequality uses
absence of an epimorphism as well as existence of a monomorphism.

The full internal chain, no-epi predicate, injection and strict-comparison
continuation is now supplied by the six linked proof units. In particular
\ref{annex:p10.lawvere-ch-chain.comparisons} excludes every parametrized
epi in the displayed directions, not just the canonical comparison map.
The qualified generalized collapse and its choice requirements remain
\ref{annex:p10.lawvere-closed-base.collapse}. None of these constructs a
specified set-theoretic relative-consistency interpretation: that still
requires an original axiom system and model, a preservation statement,
and the translation back. The introductory independence announcement
does not provide CH's positive consistency direction. Spectrum, independent
general recursion/free-algebra work and arbitrary stronger-base obligations
are not discharged by this internal forcing construction.
\end{annexitem}

\begin{annexitem}{annex:p10.lawvere-sheaves.spectrum}{thm:p10-lawvere-sheaves-spectrum}{Prime predicates, fractions, and the announced spectrum}
\sketchnote{The internal basis directness, fraction-object construction,
sheaf reflection comparisons, proof that the reflected ring is local,
and its universal local-algebra property remain to be supplied.  The
four prime predicates and monoid/fraction route are the source's
sketch, not a proof of these universal properties.}
The source route on pp.~333--334 is: form the subobject
\(X\hookrightarrow T^A\) of prime predicates; evaluate
membership to obtain the multiplicative topology basis; interpret
its incidence and inclusion relations as submonoids on the two
sides of the realization/section adjunction; form fractions there;
then reflect the presheaf fractions into the intermediate sheaf
category.  The selected multiplicative subobjects are inverted;
the article describes the resulting incidence transformation as
opposition of the membership relation.  Proving the result requires
compatible internal equivalence relations for fractions and their
universal properties, not just their object names.

The prime conditions can be imposed by finite equalizers once the
predicate/exponential and universal-quantifier maps are available.
For example multiplication gives two maps from \(T^A\) into
\(T^{A\times A}\), one evaluating at \(fg\) and the other
conjoining the evaluations at \(f\) and \(g\); their equalizer
imposes the second condition.  Evaluation at \(0\) and \(1\)
imposes the third and first.  The fourth can be expressed by the
equation \(b\land c=b\), for the predicates
\(b=[f+g\in x]\), \(c=[f\in x]\lor[g\in x]\), again
an equality of maps into \(T^{A\times A}\).
Thus finite limits impose all four universally quantified conditions;
this explains the source's finite-inverse-limit assertion without
assuming a Boolean complement operation on prime predicates.

For the ordinary classical comparison, let \(x\subseteq A\)
satisfy the four conditions and set \(\mathfrak p=A\setminus x\).
It is proper because \(1\in x\), and contains \(0\).
If \(a,b\in\mathfrak p\), the addition clause rules out
\(a+b\in x\), so \(a+b\in\mathfrak p\).  If
\(a\in\mathfrak p\) and \(r\in A\), the multiplication
equivalence rules out \(ra\in x\), hence \(ra\in\mathfrak p\).
This also gives additive inverses by taking \(r=-1\).
If \(ab\in\mathfrak p\), the equivalence says it is not the
case that both \(a,b\in x\); classical logic gives
\(a\in\mathfrak p\) or \(b\in\mathfrak p\).  Thus
\(\mathfrak p\) is a proper prime ideal.  Conversely the
complement of a proper prime ideal satisfies all four conditions:
multiplicative closure and primality give the equivalence, and
ideal closure proves the addition condition.  This converse also
uses classical complement reasoning.

For \(A=\mathbb Z\), \(\mathfrak p=(\ell)\), localization
inverts precisely the integers not divisible by \(\ell\).
A fraction with such denominator is a unit exactly when its
numerator is also not divisible by \(\ell\): one direction
inverts numerator and denominator, and the other follows from
an equation \(ac=st\) with \(s,t\) not divisible by
\(\ell\).  The nonunits form the maximal ideal generated by
\(\ell\).  This verifies the worked local-ring example but is
not the universal ringed-topos spectrum proof.
\end{annexitem}
